%
% crash_attribution.tex
%
% A Z specification of display-crash attribution: which scenes the Hub blames
% when the display ends a dial.  The state is a sequence of dials (generations
% of the connection), each with an ack ledger of the scenes written to it and
% the scenes the display acked on it; a dial is ended by the display or by the
% Hub; the death is discovered by one of several parties (the reader, the
% keepalive, an MCP tool's redial, the replicator's own send); recovery keeps a
% per-scene strike count and quarantines a scene at the threshold; after a
% charge the replicator isolates, one scene in flight at a time.
%
% This is the CONTROL, not the fixed design: a death charges nothing when it is
% discovered, but marks the dial that died; the next scene SENT is charged for it
% (docs/crash_attribution.tex is the design).  The fixed design's ruling T10 reads: when the display ends a
% dial, the suspects are exactly the scenes written to it and not acked at that
% moment; recovery charges that set once per dial, whoever discovers the death;
% a Hub-ended dial charges nothing; an empty set charges nothing.
%
% Why this model exists: the recurrence rule.  Rounds 4, 5 and 6 of the final
% review on the lux-ghhf branch each found a path on which a display death is
% blamed on the wrong scene, or on none (review-r6-silent.md F1,
% review-r6-code.md F1).  The code it governs:
% src/punt_lux/domain/hub/{batch_sender,recovery,display_connection,
% dial_generation,reply_sink,liveness_probe,replicator}.py.  ADRs: DES-102
% (attribution, isolation and quarantine), DES-093 (the Hub never controls the
% display's process), and the companion model listener_lifecycle.tex (a dial's
% generation, its ledger and its single end).
%
% Violates A2: goal  #(s,g).((s|->g):chargedOn & (s/:written(g) or s:acked(g)))
% Type-check:  fuzz docs/crash_attribution_idle_stop_buggy.tex
% Model-check (deadlock-freedom, and the disjunction of the negated goals A1 to
% A4; see Verification for the goal expression):
%   probcli docs/crash_attribution.tex -model_check -coverage -goal "<ALL>" \
%     -p MAX_OPERATIONS 100
%
% Fidelity controls (each differs from this file in the stated places only):
%   docs/crash_attribution_redial_only_buggy.tex  charge only when the
%                                                 replicator's own redial fails
%   docs/crash_attribution_idle_stop_buggy.tex    charge the next scene sent
%                                                 after a death
%   docs/crash_attribution_hub_end_charged_buggy.tex  charge a dial whoever
%                                                 ended it
%   docs/crash_attribution_double_charge_buggy.tex    charge a dial more than
%                                                 once
%
\documentclass[a4paper,10pt,fleqn]{article}
\usepackage[margin=1in]{geometry}
\usepackage{fuzz}
\usepackage[colorlinks=true,linkcolor=blue,citecolor=blue,urlcolor=blue]{hyperref}


\begin{document}

\title{Crash Attribution, charged to the next scene sent: a Control}
\author{Formal model of who is blamed when the display dies}
\date{October 2026}
\maketitle

% ============================================================
\section{Introduction}
% ============================================================

A display that crashes while it renders a scene does so a frame after it has
read the scene, so the write that carried the scene has already succeeded.  The
Hub learns of the death later, by whichever of several parties looks first: the
reader sees end of file, the keepalive's probe fails, an MCP tool redials, the
replicator's next send breaks.  The crash-loop quarantine rests on one
sentence: \emph{a scene that kills the display again and again is charged each
time, and quarantined at the threshold.}  Three review rounds in succession
found a way for that sentence to fail, each from a different side:

\begin{itemize}
  \item \textbf{Silent F1 (round~6).}\ The charge is made only where the
    replicator's own redial fails.  When the reader retires the dial and the
    keepalive or an MCP tool redials it successfully, the evidence is discarded;
    the replicator then sends every scene afresh, the display dies again, and no
    death is ever charged.  The loop runs forever.
  \item \textbf{Code F1 (round~6).}\ A write woken by the Hub's own retirement
    of the dial raises a bare socket error, which the sender charges to the
    scenes in flight; the dial was ended by the Hub and the scenes did nothing.
\end{itemize}

Round~5's finding was the same class from the other side: a deliberate stop of
the display while idle was charged to the next scene shown.  Three rounds of
one class is the recurrence signal.  The defects share a cause: \emph{the
decision whom to charge is made from the circumstances of the discovery, and
not from the dial that died.}  The ruling T10 makes the dial's own record the
only evidence.

\paragraph{The design, by responsibility.}
\begin{description}
  \item[A dial owns its ledger.]  Every dial records the scenes written to it
    ($written$) and the scenes the display acked on it ($acked$); the
    \emph{ledger} is the difference.  An ended dial keeps its ledger: it is the
    evidence of the death.
  \item[A dial has one end, with one cause.]  The display ends it (a crash, a
    deliberate stop) or the Hub does (a retire after a timeout, a drop).  The
    cause is recorded when the dial ends, by the party that ended it, and is
    never overwritten.
  \item[A death is settled once, by anyone.]  \emph{Discovery} of an ended dial
    is a step by any discoverer, and so is a redial, which discovers the death
    of the dial it replaces.  The first discovery settles the dial.  If the
    display ended it and its ledger is non-empty, that ledger is charged, one
    strike for each scene in it; otherwise nothing is charged.  Who discovers
    makes no difference to what is charged (T10).
  \item[A strike count per scene, and a threshold.]  A scene at the threshold is
    quarantined: it is dropped from the dirty set and never sent again.
  \item[Isolation narrows the suspects.]  Once any scene has a strike the
    replicator sends one scene at a time and waits for its ack before the next,
    so a later death is charged to a ledger of at most one scene.
\end{description}

\paragraph{What the model leaves out.}
The probe that detects a death while idle, the backoff between redials, and the
menu are not modelled; the display is taken to ack every scene except the
poison one, and to crash only after reading it or spontaneously.  A spontaneous
end of a dial with scenes in its ledger is indistinguishable from a poison
crash, which is the point: the ledger is the only evidence there is.

% ============================================================
\section{Constants and Free Types}
% ============================================================

Two dials: the original and its replacement, which is the least that lets a
death be discovered after a redial.  Two scenes, one of them the poison scene,
the least that lets a scene be charged for a death it did not cause.  Threshold
two: the least for which one death is not yet a quarantine, so that the loop
has a length.  Two discoverers: the replicator, and any other party.  The
traces of the controls (Section~\ref{sec:fidelity}) show that these bounds
reach both defects.

\begin{axdef}
  maxDial : \nat \\
  maxScene : \nat \\
  thr : \nat \\
  poison : \nat
\where
  maxDial = 2 \\
  maxScene = 2 \\
  thr = 2 \\
  poison = 1
\end{axdef}

\begin{zed}
  Sc == 1 \upto maxScene
\end{zed}

\begin{zed}
  Dl == 0 \upto maxDial
\end{zed}

\begin{zed}
  LV ::= lup | ldown
\end{zed}

\begin{zed}
  ENDT ::= enone | edisp | ehub
\end{zed}

\begin{zed}
  DISC ::= drepl | dother
\end{zed}

The function \emph{bump} adds a strike, and stays at the threshold.  Dial~$0$ is
a sentinel: the connection before its first dial, settled from the start.

\begin{axdef}
  bump : 0 \upto thr \fun 0 \upto thr
\where
  bump = \{ 0 \mapsto 1, 1 \mapsto 2, 2 \mapsto 2 \}
\end{axdef}

% ============================================================
\section{State}
% ============================================================

\texttt{cur} is the latest dial, \texttt{up} says whether it is live.
\texttt{ended} records the cause of each dial's end, \texttt{written} and
\texttt{acked} its ledger.  \texttt{settled} are the dials whose death a
discoverer has processed.  \texttt{dirty} are the scenes the replicator owes the
display (a redial marks every scene not quarantined).  \texttt{strikes} counts
the charges against each scene.

\texttt{chargedOn} and \texttt{times} are history, never read by a guard: the
pairs (scene, dial) that were charged, and how often each dial was charged.
The invariants are statements about them.

\begin{schema}{CA}
  cur : \nat \\
  up : LV \\
  ended : Dl \fun ENDT \\
  written : Dl \fun (\power Sc) \\
  acked : Dl \fun (\power Sc) \\
  settled : \power Dl \\
  dirty : \power Sc \\
  strikes : Sc \fun (0 \upto thr) \\
  chargedOn : Sc \rel Dl \\
  times : Dl \fun (0 \upto 2) \\
  blame : Dl
\where
  cur \leq maxDial \\
  \forall d : Dl \spot acked~d \subseteq written~d
\end{schema}

\section{Initialization}

\begin{schema}{Init}
  CA'
\where
  cur' = 0 \\
  up' = ldown \\
  \forall d : Dl \spot ended'~d = enone \\
  \forall d : Dl \spot written'~d = \emptyset \\
  \forall d : Dl \spot acked'~d = \emptyset \\
  settled' = \{ 0 \} \\
  dirty' = \emptyset \\
  \forall s : Sc \spot strikes'~s = 0 \\
  chargedOn' = \emptyset \\
  \forall d : Dl \spot times'~d = 0 \\
  blame' = 0
\end{schema}

% ============================================================
\section{Operations}
% ============================================================

The replicator's writes and the display's acks first.  \texttt{Send} writes a
scene to the live dial; in isolation the dial's ledger must be empty first (the
probe after the previous scene proved it acked).  A quarantined scene is never
sent.  \texttt{Ack} is the display's reply: it acks every scene but the poison
one.

\subsection{Send}

\begin{schema}{Send}
  \Delta CA \\
  s? : Sc
\where
  up = lup \\
  blame = 0 \\
  s? \in dirty \\
  strikes~s? < thr \\
  (\exists t : Sc \spot strikes~t > 0) \implies
    (written~cur \setminus acked~cur) = \emptyset \\
  written' = written \oplus \{ cur \mapsto written~cur \cup \{ s? \} \} \\
  dirty' = dirty \setminus \{ s? \} \\
  cur' = cur \\
  up' = up \\
  ended' = ended \\
  acked' = acked \\
  settled' = settled \\
  strikes' = strikes \\
  chargedOn' = chargedOn \\
  times' = times \\
  blame' = blame
\end{schema}

\begin{schema}{SendBlamed}
  \Delta CA \\
  s? : Sc
\where
  up = lup \\
  blame \neq 0 \\
  s? \in dirty \\
  strikes~s? < thr \\
  (\exists t : Sc \spot strikes~t > 0) \implies
    (written~cur \setminus acked~cur) = \emptyset \\
  written' = written \oplus \{ cur \mapsto written~cur \cup \{ s? \} \} \\
  dirty' = dirty \setminus \{ s? \} \\
  strikes' = strikes \oplus \{ s? \mapsto bump~(strikes~s?) \} \\
  chargedOn' = chargedOn \cup \{ s? \mapsto blame \} \\
  times' = times \oplus \{ blame \mapsto times~blame + 1 \} \\
  blame' = 0 \\
  cur' = cur \\
  up' = up \\
  ended' = ended \\
  acked' = acked \\
  settled' = settled
\end{schema}

\subsection{Ack}

\begin{schema}{Ack}
  \Delta CA \\
  s? : Sc
\where
  up = lup \\
  s? \in written~cur \setminus acked~cur \\
  s? \neq poison \\
  acked' = acked \oplus \{ cur \mapsto acked~cur \cup \{ s? \} \} \\
  cur' = cur \\
  up' = up \\
  ended' = ended \\
  written' = written \\
  settled' = settled \\
  dirty' = dirty \\
  strikes' = strikes \\
  chargedOn' = chargedOn \\
  times' = times \\
  blame' = blame
\end{schema}

\subsection{DisplayEnd and HubEnd}

The two ways a dial ends.  \texttt{DisplayEnd} is a crash or a deliberate stop,
whatever the ledger holds, empty included; \texttt{HubEnd} is a retire by the
Hub (a query timeout, a drop, a reply-skew).  Neither charges anything: a death
is charged when it is discovered, from the dial that died.

\begin{schema}{DisplayEnd}
  \Delta CA
\where
  up = lup \\
  ended' = ended \oplus \{ cur \mapsto edisp \} \\
  up' = ldown \\
  cur' = cur \\
  written' = written \\
  acked' = acked \\
  settled' = settled \\
  dirty' = dirty \\
  strikes' = strikes \\
  chargedOn' = chargedOn \\
  times' = times \\
  blame' = blame
\end{schema}

\begin{schema}{HubEnd}
  \Delta CA
\where
  up = lup \\
  ended' = ended \oplus \{ cur \mapsto ehub \} \\
  up' = ldown \\
  cur' = cur \\
  written' = written \\
  acked' = acked \\
  settled' = settled \\
  dirty' = dirty \\
  strikes' = strikes \\
  chargedOn' = chargedOn \\
  times' = times \\
  blame' = blame
\end{schema}

\subsection{Discover}

A discoverer, whoever it is, finds the current dial ended and unsettled
(\texttt{d?} is an input that no guard reads: that is T10).  If the display
ended it with a non-empty ledger, the ledger is charged and the dial settled
(\texttt{DiscoverCharge}); otherwise the dial is settled and nothing is charged
(\texttt{DiscoverQuiet}: a Hub-ended dial, or an empty ledger).  A scene that
reaches the threshold leaves the dirty set.

\begin{schema}{DiscoverBlame}
  \Delta CA \\
  d? : DISC
\where
  up = ldown \\
  cur \notin settled \\
  ended~cur = edisp \\
  settled' = settled \cup \{ cur \} \\
  blame' = cur \\
  cur' = cur \\
  up' = up \\
  ended' = ended \\
  written' = written \\
  acked' = acked \\
  dirty' = dirty \\
  strikes' = strikes \\
  chargedOn' = chargedOn \\
  times' = times
\end{schema}

\begin{schema}{DiscoverQuiet}
  \Delta CA \\
  d? : DISC
\where
  up = ldown \\
  cur \notin settled \\
  ended~cur = ehub \\
  settled' = settled \cup \{ cur \} \\
  cur' = cur \\
  up' = up \\
  ended' = ended \\
  written' = written \\
  acked' = acked \\
  dirty' = dirty \\
  strikes' = strikes \\
  chargedOn' = chargedOn \\
  times' = times \\
  blame' = blame
\end{schema}

\subsection{Redial}

A redial replaces an ended dial, and so discovers its death: the same charge,
from the same ledger, as \texttt{Discover}, made in the same step as the new
dial.  The new dial's scenes are all marked dirty except the quarantined ones.
\texttt{RedialQuiet} covers the first dial (the sentinel is settled), a dial
already settled, a Hub-ended dial and an empty ledger.

\begin{schema}{RedialBlame}
  \Delta CA \\
  d? : DISC
\where
  up = ldown \\
  cur < maxDial \\
  cur \notin settled \\
  ended~cur = edisp \\
  cur' = cur + 1 \\
  up' = lup \\
  settled' = settled \cup \{ cur \} \\
  dirty' = \{ s : Sc | strikes~s < thr \} \\
  blame' = cur \\
  ended' = ended \\
  written' = written \\
  acked' = acked \\
  strikes' = strikes \\
  chargedOn' = chargedOn \\
  times' = times
\end{schema}

\begin{schema}{RedialQuiet}
  \Delta CA \\
  d? : DISC
\where
  up = ldown \\
  cur < maxDial \\
  cur \in settled \lor ended~cur = ehub \\
  cur' = cur + 1 \\
  up' = lup \\
  settled' = settled \cup \{ cur \} \\
  strikes' = strikes \\
  dirty' = \{ s : Sc | strikes~s < thr \} \\
  chargedOn' = chargedOn \\
  times' = times \\
  ended' = ended \\
  written' = written \\
  acked' = acked \\
  blame' = blame
\end{schema}

\subsection{Halt}

\begin{schema}{Halt}
  \Xi CA
\where
  cur = maxDial \\
  up = ldown \\
  cur \in settled
\end{schema}

% ============================================================
\section{Invariants}
\label{sec:inv}
% ============================================================

Each is a negated goal: the invariant holds if its negation is unreachable
(Section~\ref{sec:verify}).  They are deliberately not conjuncts of the state
schema: in ProB's reading of Z a primed state schema is part of every
operation's guard, so an invariant written there silently disables an operation
whose result would violate it, instead of reporting the violation.

\paragraph{A1 --- a crash loop reaches quarantine, whoever discovers.}
Let a \emph{poison death} be a settled dial the display ended with the poison
scene in its ledger.  The negated goal:
\[ \# \{ g : settled \mid g > 0 \land ended~g = edisp \land
     poison \in written~g \setminus acked~g \} \ge thr
   \;\land\; strikes~poison < thr. \]
However the deaths were discovered, once $thr$ poison deaths have been settled
the poison scene is quarantined.  This is the property of the loop: it is
violated exactly when a death is settled without being charged.

\paragraph{A2 --- no scene is charged for a death it was not in.}
\[ \exists (s,g) \in chargedOn \spot
     s \notin written~g \lor s \in acked~g. \]
A scene is charged on a dial only if it was written to it and not acked on it.
In the fixed design this is a consequence of charging the ledger; its content
is that the ledger is a property of the dial, so that no circumstance of the
discovery can put another scene into it.

\paragraph{A3 --- a Hub-ended dial charges no scene.}
\[ \exists (s,g) \in chargedOn \spot ended~g = ehub. \]
The guards of \texttt{DiscoverCharge} and \texttt{RedialCharge} require the
display to have ended the dial.  Like A4 this is a tautology of the fixed
design, no operation of which charges otherwise; its content is the guard.

\paragraph{A4 --- each dial's death is charged at most once.}
\[ \exists g \in Dl \spot times~g > 1. \]
Both charging operations require the dial unsettled and settle it.  The
tautology is by the guard; \texttt{DiscoverCharge} and \texttt{RedialCharge}
race for the same dial, and the race is covered (Section~\ref{sec:verify}).

\paragraph{Deadlock-freedom.}
Every state other than the quiescent one (\texttt{Halt}) has an enabled step: a
live dial can always end, an ended one can always be discovered, and a settled
one redialed unless it is the last.

% ============================================================
\section{Fidelity Controls}
\label{sec:fidelity}
% ============================================================

\paragraph{Charge only when the replicator's own redial fails ---
\texttt{crash\_attribution\_redial\_only\_buggy.tex}.}
The code as it stands (silent F1).  \texttt{DiscoverCharge} is enabled only for
the replicator ($d? = drepl$); \texttt{RedialCharge} is deleted, so no redial
charges; \texttt{DiscoverQuiet} takes the other discoverer.  A1 is violated, in 7 steps, by
\texttt{RedialQuiet; Send(1); DisplayEnd; RedialQuiet(dother); Send(1);
DisplayEnd; DiscoverCharge(drepl)}: the first death is healed by a redial that
is not the replicator's, and its evidence is discarded; the second is charged
by the replicator, and the poison scene has one strike after two poison deaths.
(The same loop with every discoverer \texttt{dother} is longer and is also
reachable: no strike ever.)  A2, A3 and A4 hold.

\paragraph{Charge the next scene sent after a death ---
\texttt{crash\_attribution\_idle\_stop\_buggy.tex}.}
The circumstance of the discovery stands in for the dial's record.  A discovery
charges nothing but records $blame$, the dial that died; the next \texttt{Send}
carries the charge, to the scene it sends.  A2 is violated, in 4 steps, by
\texttt{RedialQuiet; DisplayEnd; RedialBlame; SendBlamed(1)}: the display was
stopped while idle (an empty ledger), and the scene shown next is charged for a
dial it was never written to.  A1 is violated too, in 7 steps, by the same
race with the poison scene in flight (\texttt{RedialQuiet; Send(1); DisplayEnd;
RedialBlame; SendBlamed(1); DisplayEnd; DiscoverBlame}: the second death is
settled with no strike).  A3 and A4 hold.

% ============================================================
\section{Verification}
\label{sec:verify}
% ============================================================

Type-checking is a \texttt{fuzz} pass:
\begin{verbatim}
  fuzz docs/crash_attribution.tex
\end{verbatim}

The model has no given sets, so no \texttt{-card}.  The check searches for a
violation of any of A1 to A4 and for deadlock; a pass is ``no counter example,
all states visited, \texttt{deadlocked:0}'':
\begin{verbatim}
  probcli docs/crash_attribution.tex -model_check -coverage -goal "<ALL>" \
    -p MAX_OPERATIONS 100
\end{verbatim}
where the goal \texttt{<ALL>} is the disjunction, in B syntax, of
\begin{verbatim}
  A1  card({g | g:settled & g>0 & ended(g)=edisp
              & 1 : written(g)-acked(g)}) >= 2 & strikes(1) < 2
  A2  #(s,g).((s|->g):chargedOn & (s/:written(g) or s:acked(g)))
  A3  #(s,g).((s|->g):chargedOn & ended(g)=ehub)
  A4  #g.(g:0..2 & times(g)>1)
\end{verbatim}
The negation of one goal, searched for with \texttt{-model\_check -nodead
-goal}, is how a control is shown to violate it.  The scenarios are reachable,
so the invariants are not vacuous (goal \emph{found} on the fixed design):
\begin{verbatim}
  strikes(1) = 2                   the poison scene is quarantined
  #g.(g:settled & g>0 & ended(g)=ehub & written(g)-acked(g) /= {})
                                   a Hub-ended dial with a full ledger, settled
  #g.(g:settled & g>0 & ended(g)=edisp & written(g)-acked(g) = {})
                                   an idle stop, settled
  #(s,g).((s|->g):chargedOn & #t.(t:1..2 & t/=s & t:written(g)))
                                   a charge made with another scene written
\end{verbatim}
Results and the bounds are recorded in \texttt{crash\_attribution\_coverage.md}.

Re-run \texttt{fuzz}, the coverage and the goals whenever
\texttt{batch\_sender.py}, \texttt{recovery.py}, \texttt{display\_connection.py},
\texttt{dial\_generation.py}, \texttt{reply\_sink.py},
\texttt{liveness\_probe.py} or \texttt{replicator.py} change who may charge a
death, from what evidence, or when.

\end{document}
